% Nouveau EGA V (cote 25): archive cote 25, site page 176 (PDF page 177).
% English translation of the checked French annotated master.
% AI-translated by GPT-6 Astra (OpenAI, ChatGPT Pro mode).
\scanpage{25/176}
\mspage{\hand{8}}

integer coefficients in the $a_i,b_{ij}$ . Then the ``content'' of this polynomial
is equal to 2 (i.e. the gcd of its coefficients is 2).
\var{\xout{That}}{}{typewriter strike-out}{ega25_25_176-v0001} Let us first prove that this content is \var{\xout{of the form $2^{\var{\deleted{r}}{}{deletion}{ega25_25_176-v0003}}$, r\illegible{1}\unsure{1}}}{}{typewriter strike-out}{ega25_25_176-v0002}\var{\inserted[typed]{a multiple of 2.}}{a multiple of 2.}{insertion}{ega25_25_176-v0004}
Indeed, this also means that \var{\deleted{\var{\xout{\unsure{(2)}}}{}{typewriter strike-out}{ega25_25_176-v0006}}}{}{deletion}{ega25_25_176-v0005} when one specializes this polynomial
in \unsure{th}\unsure{e} field $\underline{Z}\hand{/}\unsure{2|p}\underline{Z}$\hand{)} \var{\deleted{with $p\hand{\neq}2$}}{}{deletion}{ega25_25_176-v0007}, it is not identically zero, which \var{\xout{is}}{}{typewriter strike-out}{ega25_25_176-v0009}
follows from the fact that a quadratic form of \unsure{odd} degree over a field of
char 2 is always degenerate, i.e. has zero discriminant. To prove
that the content \var{\inserted{c}}{c}{insertion}{ega25_25_176-v0010} is exactly 2, it suffices to \unsure{perform} the calculation of the discriminant
\var{\inserted[typed]{of the form}}{of the form}{insertion}{ega25_25_176-v0011} standard of degree n; on the one hand it must be a multiple of c,
and on the other hand the calculation gives 2, (since it is 1 for even degree, and 2 for
degree 1 ...).

\var{\xout{One immediately concludes from this}}{}{typewriter strike-out}{ega25_25_176-v0012} One should then introduce the polynomial
\[ \overset{\hand{\sim}}{d}^{\var{\deleted{*}}{}{deletion}{ega25_25_176-v0013}}(Q) = \tfrac{1}{2}d(Q) \]
in the coefficients of Q, which again has integer coefficients, and thus
takes a well-defined value when one specializes the $a_i,b_{ij}$ in an
arbitrary ring A, i.e. when one takes a quadratic form \var{\deleted{\var{\xout{Q}}{}{typewriter strike-out}{ega25_25_176-v0015}}}{}{deletion}{ega25_25_176-v0014}\var{\inserted[typed]{q}}{q}{insertion}{ega25_25_176-v0016} of degre\unsure{e}
n with coefficients in A. We called $\overset{\hand{\sim}}{d}^{\var{\deleted{\unsure{*}}}{}{deletion}{ega25_25_176-v0017}}(Q)$ the ``\uline{corrected discriminant} \unsure{``|*}\uline{polynomial}
\uline{of the indeterminate quadratic form Q in n variables}'', and its
value with respect to the coefficients of a quadratic form q with coeff
in an arbitrary ring A will be called the ``\uline{corrected discriminant of the}
\uline{quadratic form} q''. More generally, one deduces in an essentially
trivial way from 16.5. the following statement:

\uline{Proposition} 16.6. One can, \var{\inserted[typed]{and}}{and}{insertion}{ega25_25_176-v0019} in only one way, associate with every
prescheme S equipped with a loc free Module of finite rank \uline{E} and a \var{\inserted{qu.}}{qu.}{insertion}{ega25_25_176-v0022} form
$Q \hand{\in}\var{\over{\illegible{1}}{\Gamma}}{\Gamma}{writing over}{ega25_25_176-v0023}\,\underline{\mathrm{Sym}}^2(E)$, associate a section $\overset{\hand{\sim}}{d}(Q)$ of det$(E)^{\hand{\otimes}2}$ , so as to satisfy
the following conditions:

(\unsure{a|1}) Compatibility with base change,\hand{ and functoriality for isomorphisms in \uline{E}.}

(b) If E has even rank everywhere, we have $\overset{\hand{\sim}}{d}(Q) =$ d(Q) (where d(Q)
denotes the discriminant of Q).

\mspage{\hand{176}}
\tnote{25-176-1}{The two marked parts of the retain the two doubtful letters in the source article, whose candidates permit la or un. Odd retains the a/q doubt in impair; perform retains the r/t doubt in faire; the final letter of degree retains the doubtful é. The canceled French n' has no separate English word because the live pas still gives not: the resulting statement not identically zero is retained. The canceled e in d'une and toute and the canceled d leaving e rang have no independent English effect.}
\tnote{25-176-2}{The source's pgcd is rendered gcd, char means characteristic, loc means locally, and qu. abbreviates quadratic. Its appela is rendered called, followed by the source's future will be called. The source repeats associer in Proposition 16.6; both occurrences are retained. The insertion of the form precedes standard in order to retain its variant position. The uncertain mark inside the name of the corrected discriminant polynomial remains visible, and the two nested cancellations of Q remain distinct.}
